\documentclass[12pt]{article}
\usepackage[utf8]{inputenc}
\usepackage[T1]{fontenc}
\usepackage{hyperref}
\title{Tuamotuan Language Resources for K-12 Curriculum Development}
\author{Compiled by Claude Sonnet 5.5 (claude-sonnet-5-5) for Andr\'as Kornai}
\date{}
\begin{document}
\maketitle
\section{Introduction}
Tuamotuan (Pa\okina umotu, ISO 639-3 pmt) is a Polynesian language of the Tuamotu archipelago, French Polynesia. Only a few resources appeared in search results; no modern teaching materials were found.
\section{Primary Educational Resources}
No primers or textbooks were found.
\section{Linguistic and Reference Materials}
\begin{itemize}
\item Stimson, J. F., \emph{A Dictionary of Some Tuamotuan Dialects of the Polynesian Language} (1964, 623 pp.), scanned at the Internet Archive: \url{https://archive.org/details/dictionaryofsome00stim/page/n5/mode/2up}
\item Wikipedia article on Tuamotuan: \url{https://en.wikipedia.org/wiki/Tuamotuan_language}
\end{itemize}
\section{Religious and Community Resources}
None found. A news report mentions an Acad\'emie Paumotu working to codify vocabulary and grammar of the Tuamotu areas, but no URL was seen in the results.
\section{Digital and Audio Resources}
None found.
\section{Curriculum Development Context}
\begin{itemize}
\item French Ministry of Culture languages portal, ``langue des Tuamotu'': \url{https://www.culture.gouv.fr/fr/thematiques/portail-langues-en-france/langues/langue-des-tuamotu}
\end{itemize}
Dialects vary across several linguistic areas, so any curriculum needs local community and academy involvement.
\end{document}